\section{Business Model}
\label{sec:business-model}

The monetized unit is not raw stock data. The customer pays for \textbf{risk interpretation, workflow compression, and portfolio monitoring}. The product bundles public data into a ranked daily view that answers: which stocks in my watchlist became riskier, which factor caused the change, and where should I research next?

\subsection{Packaging}

The proposed business model is subscription-based freemium:

\begin{itemize}
    \item \textbf{Free tier}: 5 tracked stocks, latest risk score, and basic overview page.
    \item \textbf{Pro tier}: NT\$199--399 per month, 20+ tracked stocks, portfolio watchlist, factor-level ranking, and historical score trends.
    \item \textbf{Community tier}: NT\$1,000--3,000 per month, shareable dashboards for finance creators, investment clubs, or paid communities.
\end{itemize}

The pricing is designed for retail affordability. A high price would force the product to compete with professional research terminals; a low-to-mid subscription lets the product compete with existing retail finance apps and newsletters.

\subsection{Revenue Drivers}

Revenue depends on three drivers:

\begin{enumerate}
    \item \textbf{Tracked-stock depth}: users are more likely to pay when the product covers enough names in their watchlist.
    \item \textbf{Trust in explanations}: each risk score must show its component drivers, not just a black-box number.
    \item \textbf{Retention through alerts}: daily or weekly risk changes create repeated usage.
\end{enumerate}

The initial go-to-market channel would be finance communities, university investment clubs, Reddit-like forums, and finance content creators. The strongest early wedge is a narrow public demo: ``20 Taiwan technology stocks ranked by explainable risk.'' If users ask to add more stocks, alerts, or portfolio import, that becomes demand evidence for the next version.
